\documentclass[11pt,letterpaper]{article}
\usepackage{cornell-notes}
\usepackage{needspace}

\newcommand{\CornellDocumentTitle}{Chapter 17: Wireless Network Security}
\title{\CornellDocumentTitle}
\author{}
\date{}
\setDocTitle{\CornellDocumentTitle}
\setDocSubtitle{Cybersecurity Cornell Notes}
\setCornellCollection{Cybersecurity Reference Notes}
\setCornellUnitType{Chapter}
\setCornellUnitNumber{17}
\setCornellUnitTitle{Wireless Network Security}
\setCornellCompanionLabel{Cybersecurity study companion}

\begin{document}
\maketitle
\begin{CornellOverviewBox}{Chapter Orientation}
\textbf{Chapter orientation.} This chapter examines wireless network security within the broader domain of overview of system and network security. Its central problem is how to reason about node, key, wireless, and routing while keeping security objectives, operating assumptions, evidence, and recovery connected. A practitioner should approach the material through protocol behavior, exposed services, trust boundaries, configuration, traffic visibility, and layered protection. Useful prerequisites include basic risk, threat, control, and systems concepts.
\end{CornellOverviewBox}
\section*{Learning Objectives}
\begin{itemize}
\item Explain the security purpose and scope of Wireless Network Security.
\item Identify the assets, actors, assumptions, and trust boundaries associated with node, key, and wireless.
\item Distinguish Chapter orientation from Cellular Networks and explain where their responsibilities intersect.
\item Analyze how failures in node, key, and wireless can affect confidentiality, integrity, availability, privacy, or safety.
\item Evaluate relevant controls through protocol behavior, exposed services, trust boundaries, configuration, traffic visibility, and layered protection.
\item Audit an implementation using network diagrams, packet evidence, rulesets, service inventories, configuration baselines, alerts, and flow records.
\item Prioritize remediation according to exploitability, consequence, control strength, and recovery readiness.
\item Verify that corrective actions change observable risk rather than documentation alone.
\end{itemize}
\section*{Cornell Cue-and-Notes}
\Needspace{8\baselineskip}
\subsection*{Chapter orientation}
\begin{CornellNotesTable}
\CornellNoteRow{What security problem does Chapter orientation address?}{\textbf{Core idea.} Treat Chapter orientation as a structured part of Wireless Network Security, with particular attention to wireless, nodes, mesh, and topology. \textbf{Analytical lens.} This topic establishes a component of the chapter's argument; connect its definitions and mechanisms to a testable security objective. For wireless, nodes, and mesh, follow traffic and state across each boundary, including malformed, unauthenticated, and partially failed conditions. \textbf{Security reasoning.} Identify what must remain protected, who or what can alter the expected state, which assumptions cross a trust boundary, and how failure changes confidentiality, integrity, availability, privacy, or safety. \textbf{Application.} The practitioner should trace traffic, validate protocol state, minimize exposure, and test prevention and detection controls. \textbf{Evidence.} Seek network diagrams, packet evidence, rulesets, service inventories, configuration baselines, alerts, and flow records.}
\end{CornellNotesTable}
\begin{CornellSummaryBox}{Topic Summary}
Chapter orientation is best remembered as the connection between wireless, nodes, mesh, and topology and the chapter's larger control objective. A defensible review states the assumption, expected behavior, evidence, failure consequence, and required response.
\end{CornellSummaryBox}
\Needspace{8\baselineskip}
\subsection*{Cellular Networks}
\begin{CornellNotesTable}
\CornellNoteRow{Which assumptions matter in Cellular Networks?}{\textbf{Core idea.} Treat Cellular Networks as a structured part of Wireless Network Security, with particular attention to wireless, nodes, base, and mobile. \textbf{Analytical lens.} This topic establishes a component of the chapter's argument; connect its definitions and mechanisms to a testable security objective. For wireless, nodes, and base, follow traffic and state across each boundary, including malformed, unauthenticated, and partially failed conditions. \textbf{Security reasoning.} Identify what must remain protected, who or what can alter the expected state, which assumptions cross a trust boundary, and how failure changes confidentiality, integrity, availability, privacy, or safety. \textbf{Application.} The practitioner should trace traffic, validate protocol state, minimize exposure, and test prevention and detection controls. \textbf{Evidence.} Seek network diagrams, packet evidence, rulesets, service inventories, configuration baselines, alerts, and flow records. Related subtopics include 802.11 Wireless LANs, Cellular Telephone Networks, and Wireless Sensor Networks.}
\end{CornellNotesTable}
\begin{CornellSummaryBox}{Topic Summary}
Cellular Networks is best remembered as the connection between wireless, nodes, base, and mobile and the chapter's larger control objective. A defensible review states the assumption, expected behavior, evidence, failure consequence, and required response.
\end{CornellSummaryBox}
\Needspace{8\baselineskip}
\subsection*{Wireless Ad Hoc Networks}
\begin{CornellNotesTable}
\CornellNoteRow{How should a reviewer evaluate Wireless Ad Hoc Networks?}{\textbf{Core idea.} Treat Wireless Ad Hoc Networks as a structured part of Wireless Network Security, with particular attention to mesh, sensor, wireless, and nodes. \textbf{Analytical lens.} This topic establishes a component of the chapter's argument; connect its definitions and mechanisms to a testable security objective. For mesh, sensor, and wireless, follow traffic and state across each boundary, including malformed, unauthenticated, and partially failed conditions. \textbf{Security reasoning.} Identify what must remain protected, who or what can alter the expected state, which assumptions cross a trust boundary, and how failure changes confidentiality, integrity, availability, privacy, or safety. \textbf{Application.} The practitioner should trace traffic, validate protocol state, minimize exposure, and test prevention and detection controls. \textbf{Evidence.} Seek network diagrams, packet evidence, rulesets, service inventories, configuration baselines, alerts, and flow records. Related subtopics include Wireless Multimedia Sensor Networks, Mesh Networks, and Internet of Things.}
\end{CornellNotesTable}
\begin{CornellSummaryBox}{Topic Summary}
Wireless Ad Hoc Networks is best remembered as the connection between mesh, sensor, wireless, and nodes and the chapter's larger control objective. A defensible review states the assumption, expected behavior, evidence, failure consequence, and required response.
\end{CornellSummaryBox}
\Needspace{8\baselineskip}
\subsection*{Security Protocols}
\begin{CornellNotesTable}
\CornellNoteRow{Why can Security Protocols fail in practice?}{\textbf{Core idea.} Treat Security Protocols as a structured part of Wireless Network Security, with particular attention to wireless, wep, access, and wired. \textbf{Analytical lens.} This topic establishes a component of the chapter's argument; connect its definitions and mechanisms to a testable security objective. For wireless, wep, and access, follow traffic and state across each boundary, including malformed, unauthenticated, and partially failed conditions. \textbf{Security reasoning.} Identify what must remain protected, who or what can alter the expected state, which assumptions cross a trust boundary, and how failure changes confidentiality, integrity, availability, privacy, or safety. \textbf{Application.} The practitioner should trace traffic, validate protocol state, minimize exposure, and test prevention and detection controls. \textbf{Evidence.} Seek network diagrams, packet evidence, rulesets, service inventories, configuration baselines, alerts, and flow records.}
\end{CornellNotesTable}
\begin{CornellSummaryBox}{Topic Summary}
Security Protocols is best remembered as the connection between wireless, wep, access, and wired and the chapter's larger control objective. A defensible review states the assumption, expected behavior, evidence, failure consequence, and required response.
\end{CornellSummaryBox}
\Needspace{8\baselineskip}
\subsection*{Wired Equivalent Privacy}
\begin{CornellNotesTable}
\CornellNoteRow{What security problem does Wired Equivalent Privacy address?}{\textbf{Core idea.} Treat Wired Equivalent Privacy as a structured part of Wireless Network Security, with particular attention to key, authentication, wpa, and message. \textbf{Analytical lens.} This topic is privacy-centered: trace collection, purpose, linkage, access, disclosure, retention, and enforceable individual protections. For key, authentication, and wpa, follow traffic and state across each boundary, including malformed, unauthenticated, and partially failed conditions. \textbf{Security reasoning.} Identify what must remain protected, who or what can alter the expected state, which assumptions cross a trust boundary, and how failure changes confidentiality, integrity, availability, privacy, or safety. \textbf{Application.} The practitioner should trace traffic, validate protocol state, minimize exposure, and test prevention and detection controls. \textbf{Evidence.} Seek network diagrams, packet evidence, rulesets, service inventories, configuration baselines, alerts, and flow records. Related subtopics include Wi-Fi Protected Access (WPA) and WPA2, Wpa, SPINS: Security Protocols for Sensor Networks, and Secure Network Encryption Protocol.}
\end{CornellNotesTable}
\begin{CornellSummaryBox}{Topic Summary}
Wired Equivalent Privacy is best remembered as the connection between key, authentication, wpa, and message and the chapter's larger control objective. A defensible review states the assumption, expected behavior, evidence, failure consequence, and required response.
\end{CornellSummaryBox}
\Needspace{8\baselineskip}
\subsection*{Secure Routing}
\begin{CornellNotesTable}
\CornellNoteRow{Which assumptions matter in Secure Routing?}{\textbf{Core idea.} Treat Secure Routing as a structured part of Wireless Network Security, with particular attention to routing, node, route, and sequence. \textbf{Analytical lens.} This topic is control-centered: map each safeguard to a specific failure mode, then test independence, coverage, and bypass paths. For routing, node, and route, follow traffic and state across each boundary, including malformed, unauthenticated, and partially failed conditions. \textbf{Security reasoning.} Identify what must remain protected, who or what can alter the expected state, which assumptions cross a trust boundary, and how failure changes confidentiality, integrity, availability, privacy, or safety. \textbf{Application.} The practitioner should trace traffic, validate protocol state, minimize exposure, and test prevention and detection controls. \textbf{Evidence.} Seek network diagrams, packet evidence, rulesets, service inventories, configuration baselines, alerts, and flow records. Related subtopics include Secure Efficient Ad Hoc Distance, and Ariadne.}
\end{CornellNotesTable}
\begin{CornellSummaryBox}{Topic Summary}
Secure Routing is best remembered as the connection between routing, node, route, and sequence and the chapter's larger control objective. A defensible review states the assumption, expected behavior, evidence, failure consequence, and required response.
\end{CornellSummaryBox}
\Needspace{8\baselineskip}
\subsection*{Authenticated Routing For Ad Hoc Networks}
\begin{CornellNotesTable}
\CornellNoteRow{How should a reviewer evaluate Authenticated Routing For Ad Hoc Networks?}{\textbf{Core idea.} Treat Authenticated Routing For Ad Hoc Networks as a structured part of Wireless Network Security, with particular attention to node, rep, rdp, and signature. \textbf{Analytical lens.} This topic establishes a component of the chapter's argument; connect its definitions and mechanisms to a testable security objective. For node, rep, and rdp, follow traffic and state across each boundary, including malformed, unauthenticated, and partially failed conditions. \textbf{Security reasoning.} Identify what must remain protected, who or what can alter the expected state, which assumptions cross a trust boundary, and how failure changes confidentiality, integrity, availability, privacy, or safety. \textbf{Application.} The practitioner should trace traffic, validate protocol state, minimize exposure, and test prevention and detection controls. \textbf{Evidence.} Seek network diagrams, packet evidence, rulesets, service inventories, configuration baselines, alerts, and flow records.}
\end{CornellNotesTable}
\begin{CornellSummaryBox}{Topic Summary}
Authenticated Routing For Ad Hoc Networks is best remembered as the connection between node, rep, rdp, and signature and the chapter's larger control objective. A defensible review states the assumption, expected behavior, evidence, failure consequence, and required response.
\end{CornellSummaryBox}
\Needspace{8\baselineskip}
\subsection*{Secure Link State Routing Protocol}
\begin{CornellNotesTable}
\CornellNoteRow{Why can Secure Link State Routing Protocol fail in practice?}{\textbf{Core idea.} Treat Secure Link State Routing Protocol as a structured part of Wireless Network Security, with particular attention to node, certificate, lsu, and key. \textbf{Analytical lens.} This topic is control-centered: map each safeguard to a specific failure mode, then test independence, coverage, and bypass paths. For node, certificate, and lsu, follow traffic and state across each boundary, including malformed, unauthenticated, and partially failed conditions. \textbf{Security reasoning.} Identify what must remain protected, who or what can alter the expected state, which assumptions cross a trust boundary, and how failure changes confidentiality, integrity, availability, privacy, or safety. \textbf{Application.} The practitioner should trace traffic, validate protocol state, minimize exposure, and test prevention and detection controls. \textbf{Evidence.} Seek network diagrams, packet evidence, rulesets, service inventories, configuration baselines, alerts, and flow records. Related subtopics include Bootstrapping in Wireless Ad Hoc Networks.}
\end{CornellNotesTable}
\begin{CornellSummaryBox}{Topic Summary}
Secure Link State Routing Protocol is best remembered as the connection between node, certificate, lsu, and key and the chapter's larger control objective. A defensible review states the assumption, expected behavior, evidence, failure consequence, and required response.
\end{CornellSummaryBox}
\Needspace{8\baselineskip}
\subsection*{Key Establishment}
\begin{CornellNotesTable}
\CornellNoteRow{What security problem does Key Establishment address?}{\textbf{Core idea.} Treat Key Establishment as a structured part of Wireless Network Security, with particular attention to nodes, wireless, sensor, and keys. \textbf{Analytical lens.} This topic is cryptography-centered: state the intended property and validate algorithms, parameters, keys, protocol binding, and lifecycle controls. For nodes, wireless, and sensor, follow traffic and state across each boundary, including malformed, unauthenticated, and partially failed conditions. \textbf{Security reasoning.} Identify what must remain protected, who or what can alter the expected state, which assumptions cross a trust boundary, and how failure changes confidentiality, integrity, availability, privacy, or safety. \textbf{Application.} The practitioner should trace traffic, validate protocol state, minimize exposure, and test prevention and detection controls. \textbf{Evidence.} Seek network diagrams, packet evidence, rulesets, service inventories, configuration baselines, alerts, and flow records. Related subtopics include Bootstrapping, Bootstrapping in Wireless Sensor Networks, Contributory Schemes, and Key Management.}
\end{CornellNotesTable}
\begin{CornellSummaryBox}{Topic Summary}
Key Establishment is best remembered as the connection between nodes, wireless, sensor, and keys and the chapter's larger control objective. A defensible review states the assumption, expected behavior, evidence, failure consequence, and required response.
\end{CornellSummaryBox}
\Needspace{8\baselineskip}
\subsection*{Ingemarsson, Tang, And Wong}
\begin{CornellNotesTable}
\CornellNoteRow{Which assumptions matter in Ingemarsson, Tang, And Wong?}{\textbf{Core idea.} Treat Ingemarsson, Tang, And Wong as a structured part of Wireless Network Security, with particular attention to node, key, group, and value. \textbf{Analytical lens.} This topic establishes a component of the chapter's argument; connect its definitions and mechanisms to a testable security objective. For node, key, and group, follow traffic and state across each boundary, including malformed, unauthenticated, and partially failed conditions. \textbf{Security reasoning.} Identify what must remain protected, who or what can alter the expected state, which assumptions cross a trust boundary, and how failure changes confidentiality, integrity, availability, privacy, or safety. \textbf{Application.} The practitioner should trace traffic, validate protocol state, minimize exposure, and test prevention and detection controls. \textbf{Evidence.} Seek network diagrams, packet evidence, rulesets, service inventories, configuration baselines, alerts, and flow records. Related subtopics include Distributed Schemes, Hypercube and Octopus, Partially Distributed Threshold Certificate, and Authority Scheme.}
\end{CornellNotesTable}
\begin{CornellSummaryBox}{Topic Summary}
Ingemarsson, Tang, And Wong is best remembered as the connection between node, key, group, and value and the chapter's larger control objective. A defensible review states the assumption, expected behavior, evidence, failure consequence, and required response.
\end{CornellSummaryBox}
\Needspace{8\baselineskip}
\subsection*{Management Countermeasures}
\begin{CornellNotesTable}
\CornellNoteRow{How should a reviewer evaluate Management Countermeasures?}{\textbf{Core idea.} Treat Management Countermeasures as a structured part of Wireless Network Security, with particular attention to key, wireless, management, and phase. \textbf{Analytical lens.} This topic is control-centered: map each safeguard to a specific failure mode, then test independence, coverage, and bypass paths. For key, wireless, and management, follow traffic and state across each boundary, including malformed, unauthenticated, and partially failed conditions. \textbf{Security reasoning.} Identify what must remain protected, who or what can alter the expected state, which assumptions cross a trust boundary, and how failure changes confidentiality, integrity, availability, privacy, or safety. \textbf{Application.} The practitioner should trace traffic, validate protocol state, minimize exposure, and test prevention and detection controls. \textbf{Evidence.} Seek network diagrams, packet evidence, rulesets, service inventories, configuration baselines, alerts, and flow records.}
\end{CornellNotesTable}
\begin{CornellSummaryBox}{Topic Summary}
Management Countermeasures is best remembered as the connection between key, wireless, management, and phase and the chapter's larger control objective. A defensible review states the assumption, expected behavior, evidence, failure consequence, and required response.
\end{CornellSummaryBox}
\begin{CornellWarningBox}{Historical Context}
Some products, protocols, examples, threat observations, or legal references in this chapter are historically bounded. Retain the underlying threat model and control objective, but verify present applicability before treating a named implementation as current guidance.
\end{CornellWarningBox}
\begin{CornellOverviewBox}{Defensive Guidance}
Apply the chapter by making assets, authority, trust, expected behavior, and failure response explicit. In practice, trace traffic, validate protocol state, minimize exposure, and test prevention and detection controls. A control is not fully supported until its operation, monitoring, ownership, and recovery behavior are demonstrated with evidence.
\end{CornellOverviewBox}
\section*{Threat-and-Control Map}
\begingroup\scriptsize\setlength{\tabcolsep}{2pt}
\begin{longtable}{>{\raggedright\arraybackslash}p{0.135\textwidth}>{\raggedright\arraybackslash}p{0.145\textwidth}>{\raggedright\arraybackslash}p{0.155\textwidth}>{\raggedright\arraybackslash}p{0.145\textwidth}>{\raggedright\arraybackslash}p{0.16\textwidth}>{\raggedright\arraybackslash}p{0.17\textwidth}}
\toprule\rowcolor{Navy}\color{white}\textbf{Asset} & \color{white}\textbf{Threat / Failure} & \color{white}\textbf{Vulnerability} & \color{white}\textbf{Impact} & \color{white}\textbf{Detection / Evidence} & \color{white}\textbf{Control}\\\midrule
\endfirsthead
\toprule\rowcolor{Navy}\color{white}\textbf{Asset} & \color{white}\textbf{Threat / Failure} & \color{white}\textbf{Vulnerability} & \color{white}\textbf{Impact} & \color{white}\textbf{Detection / Evidence} & \color{white}\textbf{Control}\\\midrule
\endhead
Network services, communications, and connected hosts & Unauthorized access, interception, manipulation, or disruption & Exposed services, weak protocol assumptions, or unsafe configuration & Loss of confidentiality, integrity, availability, or control & Packet, flow, host, and alert telemetry correlated to expected behavior & Segmentation, hardened configuration, authentication, filtering, and monitoring\\\midrule
Assumptions surrounding node & Misuse or unexpected state transition & Trust in key is not validated & A local weakness becomes a broader compromise path & Review evidence for wireless and test negative paths & Make trust explicit, constrain authority, and fail safely\\\midrule
Recovery capability and decision evidence & Control failure remains unnoticed or cannot be reversed & Monitoring, ownership, or restoration has not been exercised & Longer exposure and slower, less reliable recovery & Alert-to-response exercises and restoration tests & Assigned ownership, actionable telemetry, preserved evidence, and rehearsed recovery\\\midrule
\bottomrule
\end{longtable}
\endgroup
\section*{Practical Security Checklist}
\begin{CornellChecklist}
\item Define the in-scope assets, users, services, data, and operational dependencies for Wireless Network Security.
\item Document the expected behavior and security objective of node.
\item Map identities, privileges, trust boundaries, and externally controlled inputs associated with node and key.
\item Record the threat or failure being addressed, its prerequisites, likely path, and credible consequences.
\item Inspect network diagrams, packet evidence, rulesets, service inventories, configuration baselines, alerts, and flow records.
\item Test both the intended path and negative paths, including invalid state, partial failure, excessive privilege, and unavailable dependencies.
\item Confirm preventive controls are complemented by actionable detection and a named response owner.
\item Exercise containment, restoration, and wireless; record actual recovery behavior and unresolved dependencies.
\item Validate exceptions, compensating controls, expiration dates, and accountable approval.
\item Preserve reproducible evidence and tie each finding to affected assets, risk, remediation, and verification criteria.
\end{CornellChecklist}
\begin{CornellSummaryBox}{Chapter Synthesis}
The chapter's principal lesson is that wireless network security must be evaluated as an operating security system rather than an isolated product or statement. The most important failure patterns involve node, key, wireless, and routing, especially when ownership, boundaries, telemetry, or recovery remain implicit. A reviewer should connect architecture and behavior to network diagrams, packet evidence, rulesets, service inventories, configuration baselines, alerts, and flow records, prioritize findings by reachable consequence and control independence, and verify remediation through observable change. Within Overview of System and Network Security, this chapter contributes a reusable method: state the objective, expose assumptions, test failure, preserve evidence, and learn from operation.
\end{CornellSummaryBox}
\section*{Self-Test Questions}
\begin{CornellRecallList}
\item What is the central security objective of Wireless Network Security?
\item How does Chapter orientation contribute to the chapter's broader security argument?
\item Distinguish Cellular Networks from Wireless Ad Hoc Networks.
\item Which trust assumptions should be made explicit when evaluating node?
\item How could a weakness involving key affect confidentiality, integrity, availability, privacy, or safety?
\item What evidence would demonstrate that controls surrounding wireless operate as intended?
\item Why should prevention, detection, response, and recovery be evaluated together in this domain?
\item Scenario: A reachable service accepts traffic across a boundary but its assumptions and monitoring are undocumented. What should the reviewer examine first, and why?
\item Scenario: A control associated with routing passes a configuration check but fails during an operational exercise. How should the finding be interpreted and prioritized?
\item How should a practitioner distinguish enduring principles in this chapter from time-sensitive tools, examples, or implementation details?
\end{CornellRecallList}
\Needspace{8\baselineskip}
\section*{Answer Key}
\begin{enumerate}
\item The objective is to protect the relevant assets by understanding protocol behavior, exposed services, trust boundaries, configuration, traffic visibility, and layered protection and by connecting controls to measurable risk reduction.
\item Chapter orientation provides one part of the causal chain: it should identify expected behavior, dependencies, failure modes, and the evidence needed to justify confidence.
\item Cellular Networks and Wireless Ad Hoc Networks address different portions of the problem. A strong comparison states their scope, assumptions, enforcement points, evidence, and how a failure in one affects the other.
\item Reviewers should identify who controls node, what inputs or dependencies it trusts, where privilege changes, what happens during partial failure, and how those assumptions are tested.
\item A weakness in key can propagate across trust boundaries. The actual impact depends on reachable assets, granted authority, control independence, visibility, and recovery capability.
\item Useful evidence includes network diagrams, packet evidence, rulesets, service inventories, configuration baselines, alerts, and flow records; configuration alone is insufficient when operation, monitoring, or recovery is part of the claim.
\item Prevention reduces probability, detection limits dwell time, response limits spread, and recovery restores objectives. Omitting one layer creates an unexamined failure mode.
\item Begin by establishing scope, ownership, expected behavior, and the violated assumption. Then trace traffic, validate protocol state, minimize exposure, and test prevention and detection controls, preserving evidence for prioritization and follow-up.
\item Treat the exercise as stronger evidence than a static check. Prioritize according to reachable impact, likelihood of real operating conditions, missing compensating controls, and recovery difficulty.
\item Retain the principle, threat model, and control objective; separately label technologies, legal references, product examples, and threat observations whose applicability depends on time or environment.
\end{enumerate}
\section*{Key-Term Recap}
\begin{longtable}{>{\raggedright\arraybackslash}p{0.27\textwidth}>{\raggedright\arraybackslash}p{0.66\textwidth}}
\toprule\rowcolor{Navy}\color{white}\textbf{Term} & \color{white}\textbf{Meaning}\\\midrule
\endfirsthead
\toprule\rowcolor{Navy}\color{white}\textbf{Term} & \color{white}\textbf{Meaning}\\\midrule
\endhead
\textbf{Packet} & A formatted unit of network-layer communication containing control fields and payload data. \\\midrule
\textbf{Authentication} & The process of establishing confidence that an asserted identity is genuine. \\\midrule
\textbf{Certificate} & A signed data structure that binds identity information to a public key under a trust model. \\\midrule
\textbf{Cipher} & An algorithm that transforms data using a key to provide a defined cryptographic security property. \\\midrule
\textbf{Wireless Sensor Network} & A network of sensing devices whose resource constraints and unattended deployment shape its threat model. \\\midrule
\textbf{Risk} & The effect of uncertainty on objectives, commonly evaluated through likelihood, impact, exposure, and context. \\\midrule
\textbf{Access Control} & Rules and enforcement mechanisms that restrict actions on resources to authorized subjects. \\\midrule
\textbf{Integrity} & The property that information and systems remain accurate, complete, and protected from unauthorized modification. \\\midrule
\textbf{Aes} & A standardized symmetric block cipher used to protect data when implemented with an appropriate mode and key lifecycle. \\\midrule
\textbf{Encryption} & A keyed transformation of plaintext into ciphertext to protect confidentiality. \\\midrule
\bottomrule
\end{longtable}
\begin{CornellExamBox}{One-Sentence Takeaway}
Treat wireless network security as a verifiable chain from assets and assumptions through controls, evidence, response, and recovery.
\end{CornellExamBox}
\end{document}
